\section*{Hoofdstuk \ref{ch:b3:many-electron-atoms} --- Meerelektronatomen en termsymbolen}

\begin{solution}{exo:b3:many-electron-atoms:1}
\[
\Psi = \frac{1}{\sqrt6}\begin{vmatrix}1s\alpha(1) & 1s\beta(1) & 2s\alpha(1)\\
1s\alpha(2) & 1s\beta(2) & 2s\alpha(2)\\ 1s\alpha(3) & 1s\beta(3) & 2s\alpha(3)
\end{vmatrix}.
\]
$1s$ biedt maar twee spinorbitalen, $1s\alpha$ en $1s\beta$; een derde
elektron zou er een van herhalen, waardoor twee kolommen gelijk worden: de
determinant is nul.
\end{solution}

\begin{solution}{exo:b3:many-electron-atoms:2}
$\binom62 = 15$, $\binom63 = 20$, $\binom{10}{2} = 45$. $p^3$: ${}^4$S (4) +
${}^2$D (10) + ${}^2$P (6) = 20. $d^2$: ${}^3$F (21) + ${}^3$P (9) + ${}^1$G (9) +
${}^1$D (5) + ${}^1$S (1) = 45.
\end{solution}

\begin{solution}{exo:b3:many-electron-atoms:3}
N ($2p^3$, halfvol): \termsym{4}{S}{3/2}. O ($2p^4$, meer dan halfvol):
\termsym{3}{P}{2}. F ($2p^5$): \termsym{2}{P}{3/2}. \ce{Ti^2+} ($3d^2$):
\termsym{3}{F}{2}. \ce{Fe^3+} ($3d^5$, halfvol, alle spins parallel, $L = 0$):
\termsym{6}{S}{5/2}.
\end{solution}

\begin{solution}{exo:b3:many-electron-atoms:4}
${}^2$D: $J = \frac52$ (6 toestanden) en $\frac32$ (4); $6 + 4 = 10 = 5 \times 2$.
${}^3$P: $J = 2$ (5), 1 (3), 0 (1); $5 + 3 + 1 = 9 = 3 \times 3$.
\end{solution}

\begin{solution}{exo:b3:many-electron-atoms:5}
Intervallen $16.4$ ($J = 1 - 0$) en $27.0$ ($2 - 1$): verhouding 1.65 in plaats van 2.
$A = 16.4/1 = \qty{16.4}{cm^{-1}}$ uit het eerste, $27.0/2 = \qty{13.5}{cm^{-1}}$
uit het tweede. Eén enkele $A$ past niet: de niveaus zijn licht gemengd met
die van andere termen (${}^1$D, ${}^1$S), en \omterm{def:b3:many-electron-atoms:russell-saunders}{russell-saunderskoppeling} is een
benadering.
\end{solution}

\begin{solution}{exo:b3:many-electron-atoms:6}
$\lambda = 10^7/\tilde\nu$: \qty{589.76}{nm} (D$_1$) en \qty{589.16}{nm}
(D$_2$). Afstand $\qty{17.20}{cm^{-1}} = \qty{2.13}{meV}$. $E(\frac32) -
E(\frac12) = \frac32A$, dus $A = \qty{11.47}{cm^{-1}}$.
\end{solution}

\begin{solution}{exo:b3:many-electron-atoms:7}
$3p \to 3s$: toegestaan ($\Delta l = -1$). $3d \to 3s$: verboden ($\Delta l = -2$).
$3d \to 3p$: toegestaan. $4s \to 3s$: verboden ($\Delta l = 0$, geen
verandering van pariteit). $4s \to 3p$: toegestaan.
\end{solution}

\begin{solution}{exo:b3:many-electron-atoms:8}
Gemiddelde van ${}^3$P: $(5 \times 169086.76 + 3 \times 169086.84 + 169087.83)/9 =
\qty{169086.91}{cm^{-1}}$. Afstand tot ${}^1$P: \qty{2047.98}{cm^{-1}} $=
\qty{0.254}{eV}$, dus $K(1s,2p) = \qty{0.127}{eV}$, drie keer kleiner dan
$K(1s,2s)$. De \omterm{def:b3:many-electron-atoms:exchange}{uitwisselingsintegraal} meet de overlap van de dichtheden van
de twee orbitalen; de $2s$-orbitaal dringt door in het gebied van $1s$ (ze
heeft dichtheid dicht bij de kern), de $2p$-orbitaal is nul in de kern en
overlapt minder.
\end{solution}

\begin{solution}{exo:b3:many-electron-atoms:9}
De grootste $M_L = 4$ vereist beide elektronen in $m_l = 2$, met
tegengestelde spins: een singulet, ${}^1$G (9 toestanden). De volgende,
$M_L = 3$, komt van $m_l = 2, 1$, met
$M_S = 1$ mogelijk: ${}^3$F (21). Er blijft $M_L = 2$ over met alleen $M_S = 0$:
${}^1$D (5). Er blijft $M_L = 1$ over met $M_S = 1$: ${}^3$P (9). Eén toestand blijft over met
$M_L = M_S = 0$: ${}^1$S. Totaal 45.
\end{solution}

\begin{solution}{exo:b3:many-electron-atoms:10}
Voor $\mathbf r_1 = \mathbf r_2 = \mathbf r$ wordt de tripletfunctie
$\frac{1}{\sqrt2}[1s(\mathbf r)2s(\mathbf r) - 2s(\mathbf r)1s(\mathbf r)] = 0$; de
singuletfunctie is $\sqrt2\,1s(\mathbf r)2s(\mathbf r) \ne 0$. In de \omterm{def:b3:many-electron-atoms:singlet-triplet}{triplettoestand}
worden de elektronen nooit in hetzelfde punt aangetroffen en liggen ze
gemiddeld verder uit elkaar, zodat hun afstoting kleiner is: $E_+ - E_- = 2K > 0$.
\end{solution}

\begin{solution}{exo:b3:many-electron-atoms:11}
$d^8$ is meer dan halfvol: de koppeling is omgekeerd, $J = L + S = 4$
ligt het laagst. Intervallen: $E(3) - E(4) = 1360.7$ en $E(2) - E(3) = 908.9$;
verhouding 1.50, tegenover $4/3 = 1.33$ volgens Landé. $|A| = 1360.7/4 = \qty{340}{cm^{-1}}$
($A < 0$). De \omterm{def:b3:many-electron-atoms:spin-orbit}{spin-baankoppeling} is in dit zwaardere ion veel groter dan in
koolstof (\qty{16}{cm^{-1}}).
\end{solution}

\begin{solution}{exo:b3:many-electron-atoms:12}
$\sum_J(2J+1)[J(J+1) - L(L+1) - S(S+1)] = 0$ (de $(2L+1)(2S+1)$ toestanden
kunnen in beide bases geteld worden, en $\sum M_LM_S = 0$ over de
ongekoppelde). Voor ${}^3$P: $1(0 - 4) + 3(2 - 4) + 5(6 - 4) = -4 - 6 + 10 = 0$.
Het gewogen gemiddelde van koolstof is $(0 + 3 \times 16.42 + 5 \times 43.41)/9 = \qty{29.6}{cm^{-1}}$:
de energie die de term zonder \omterm{def:b3:many-electron-atoms:spin-orbit}{spin-baankoppeling} zou hebben.
\end{solution}

\begin{solution}{pb:b3:many-electron-atoms:1}
\textbf{1.} $1s\alpha$, $1s\beta$, $2s\alpha$, $2s\beta$.
\textbf{2.} $|1s\alpha\,2s\alpha|$, $|1s\alpha\,2s\beta|$, $|1s\beta\,2s\alpha|$,
$|1s\beta\,2s\beta|$.
\textbf{3.} $|1s\alpha\,2s\alpha| = \frac{1}{\sqrt2}[1s(1)2s(2) - 2s(1)1s(2)]\,
\alpha(1)\alpha(2)$, door de $2\times2$-determinant uit te werken.
\textbf{4.} Hun som is $\frac{1}{\sqrt2}[1s(1)2s(2) - 2s(1)1s(2)]\cdot
\frac{1}{\sqrt2}[\alpha(1)\beta(2) + \beta(1)\alpha(2)]$ (maal $\sqrt2$, daarna
genormeerd); hun verschil geeft $\frac{1}{\sqrt2}[1s(1)2s(2) +
2s(1)1s(2)]\cdot\frac{1}{\sqrt2}[\alpha(1)\beta(2) - \beta(1)\alpha(2)]$.
\textbf{5.} Symmetrische ruimtefunctie met de antisymmetrische spinfunctie
($S = 0$); antisymmetrische ruimtefunctie met de drie symmetrische
spinfuncties ($S = 1$).
\textbf{6.} $2S + 1 = 1$ en 3: één en drie toestanden met dezelfde energie.
\textbf{7.} $[\ce{Ar}]\,3d^2$; $\binom{10}{2} = 45$ determinanten.
\textbf{8.} Maximale $S = 1$, maximale $L = 2 + 1 = 3$, minder dan halfvol:
\termsym{3}{F}{2}, het niveau bij 0 in de gegevens.
\textbf{9.} $21 + 9 + 9 + 5 + 1 = 45$.
\textbf{10.} Normaal ($J = 2$ het laagst). Intervallen 184.9 en 235.5:
verhouding 1.27 tegenover $4/3 = 1.33$; de regel klopt binnen 5\,\%.
\textbf{11.} ${}^3$F: $(5 \times 0 + 7 \times 184.9 + 9 \times 420.4)/21 =
\qty{241.8}{cm^{-1}}$; ${}^3$P: $(10538.4 + 3 \times 10603.6 + 5 \times
10721.2)/9 = \qty{10661.7}{cm^{-1}}$.
\textbf{12.} $15B = \qty{10419.9}{cm^{-1}}$, $B = \qty{695}{cm^{-1}}$.
\textbf{13.} $\Delta S = 0$.
\textbf{14.} $\Delta S = 1$ is verboden, en $1s2s \to 1s^2$ heeft $\Delta l = 0$
(geen verandering van pariteit; ${}^3$S$_1 \to {}^1$S$_0$ heeft bovendien
$\Delta L = 0$ tussen twee S-toestanden). De toestand is metastabiel: hij
leeft veel langer dan gewone aangeslagen toestanden.
\textbf{15.} $169086.91 - 159855.97 = \qty{9230.94}{cm^{-1}}$: \qty{1083.3}{nm}, in
het nabije infrarood.
\textbf{16.} $171134.89 - 166277.44 = \qty{4857.45}{cm^{-1}}$: \qty{2058.7}{nm}.
\textbf{17.} $\qty{171134.89}{cm^{-1}}$: \qty{58.43}{nm}, in het
vacuümultraviolet (geabsorbeerd door lucht).
\textbf{18.} Elke familie heeft haar eigen lijnen en haar eigen laagste
toestand, en de twee zijn nooit door licht met elkaar verbonden: ze
gedragen zich als twee gassen met twee spectra.
\textbf{19.} $E_\pm = E_{1s} + E_{2s} + J \pm K$ (singulet $+$, triplet $-$).
\textbf{20.} De triplet: zijn ruimtelijke functie is nul wanneer de
elektronen elkaar ontmoeten, zodat ze elkaar minder afstoten.
\textbf{21.} \qty{6421.47}{cm^{-1}}, $6421.47/8065.54 = \qty{0.796}{eV}$.
\textbf{22.} Afstand $\qty{2047.98}{cm^{-1}} = \qty{0.254}{eV}$, dus $K(1s,2p) =
\qty{0.127}{eV}$.
\textbf{23.} $2s$ dringt door in het gebied van $1s$ en overlapt het meer
dan $2p$.
\textbf{24.} Ja: in beide configuraties ligt de triplet (grotere $S$) het laagst.
\textbf{25.} $K = \frac12 \times \qty{0.796}{eV}$: \textbf{de
\omterm{def:b3:many-electron-atoms:exchange}{uitwisselingsintegraal} $K(1s,2s)$ van helium is \qty{0.398}{eV}}, gemeten
als de helft van de singulet-tripletafstand.
\end{solution}
